% Introduction to the bishop vs bishop (same colour) chapter, for the 1500-1599 book only.
% There is no shared texts/bishop-vs-bishop-same-colour.tex yet; if one is written, \input it at the end.
% Words defined further down in this chapter are defined here as well, so that the links work from
% the start; the first definition in the book wins.

\section*{What to look for}
With bishops of the same colour, both bishops fight for the same squares, so the bishop can attack
or block the enemy bishop directly. The extra pawn often does not win, and the defender has more
resources than it seems. Most mistakes in this chapter came from a few ideas.

\begin{itemize}
\item The defender gives up his bishop for the last pawn. Without pawns, a bishop up cannot win
(\term{removing-the-pawns}).\index{Exchanging or removing the pawns} Even one
pawn left over may not be enough: bishop and rook's pawn draw if the bishop is the
\vocab[bishop endings]{wrong-bishop}{wrong bishop}{A bishop that does not control the promotion
square of a rook's pawn. With such a bishop and a rook's pawn, the stronger side cannot drive the
defending king out of the corner, and the game is a draw.}.\index{Theoretical draws!Wrong bishop and rook's pawn against king}
Both sides missed this again and again: \xref{Diagram}{stefansson-opasiak-74w},
\xref{Diagram}{christiansen-edwardsson-68b} and \xref{Diagram}{boitumelosethabathaba-magnekouokamsylviane-72b}.

\item The attacker's typical weapon is \term{interference}:\index{Bishop!Interference} his bishop
steps onto the defending bishop's diagonal and cuts it off from the queening square. The defender
answers with the
\vocab[bishop endings]{king-after-king}{king after king}{The defending king follows the attacking
king and stands opposite it, so that the attacker's bishop cannot interfere with the defending
bishop's diagonal.} defence\index{Bishop!King after king} (\xref{Diagram}{arcecano-grigorchuk-65w},
\xref{Diagram}{debankurbanerjee-mohit-76b} and \xref{Diagram}{debankurbanerjee-mohit-69w}), and the
attacker tries
\vocab[bishop endings]{drive-off-the-bishop}{driving off the bishop}{Attacking the enemy bishop, or
offering an exchange, so that it must leave the diagonal on which it does its job.}\index{Bishop!Drive off the bishop}
(\xref{Diagram}{aliyev-baghirov-54w}).

\item The defending bishop works best \termas{one-diagonal}{along one
diagonal}.\index{Bishop!One-diagonal principle} Together with the pawns it can set up a
\term{barrier}\index{Barrier} that the enemy king cannot cross. Put your pawns on the
\vocab[bishop endings]{pawns-other-colour}{other colour than the bishop}{Put your pawns on squares
of the colour your bishop does not move on. Then bishop and pawns together control squares of both
colours, the bishop does not have to guard its own pawns, and the pawns do not block its
diagonals.}:\index{Bishop!Pawns on the other colour than the bishop} then bishop and pawns guard
squares of both colours (\xref{Diagram}{butti-torrescaballero-54w},
\xref{Diagram}{aragonda-rekhviashvili-44w} and \xref{Diagram}{cortesdiaz-deleonvillarreal-49w}).

\item Many of these positions are a \term{fortress}:\index{Fortress} the defender only has to keep
it, and the attacker must find a way in (\xref{Diagram}{ilbay-daloglu-53w},
\xref{Diagram}{fragiadakis-sempekopoulos-40b} and \xref{Diagram}{khafizov-nadrshin-71w}). Use the
\term{process-of-elimination}\index{Practical play!Process of elimination} to find the one move
that keeps it.

\item The attacker wins with patience. \termas{do-not-rush}{Do not rush},\index{Strategy!Do not rush}
\termas{do-not-weaken}{do not weaken yourself},\index{Strategy!Do not weaken yourself} and first take
away the opponent's only resource with \term{prophylaxis}\index{Strategy!Prophylaxis}
(\xref{Diagram}{isak-cam-44b}, \xref{Diagram}{liapizenebisi-kalyva-101w} and
\xref{Diagram}{teshima-sarda-60b}). Often
\vocab[bishop endings]{bishop-escorts}{the bishop escorts the pawn}{The bishop supports the pawn's
advance by guarding the squares in front of it, so that the king is free for other
tasks.}\index{Bishop!Escorting the pawn} while the king does other work
(\xref{Diagram}{kwiek-kuznik-39b}).

\item When both sides have passed pawns, use \term{counting}\index{Practical play!Counting} before
you act. In a pawn race every tempo counts:
\vocab[pawn races]{do-not-lose-tempo}{do not lose a tempo}{In a pawn race every tempo counts: a single
move that does not bring a pawn closer to its queening square, or does not stop the enemy pawn,
can decide the race.},\index{Pawn races!Do not lose a tempo} and it also
matters who queens with check\index{Pawn races!Queening with check} (\xref{Diagram}{jiang-han-57b},
\xref{Diagram}{fazla-ersoy-51w} and \xref{Diagram}{moretto-cardoso-49b}).
\end{itemize}
